% SEGMENT OLP-0595-S01
\documentclass[../../../include/open-logic-section]{subfiles}

\begin{document}

\olfileid{sth}{choice}{hartogs}
\olsection{ஒப்பிடத்தக்க தன்மையும் ஹார்டாக்ஸ் துணைத்தேற்றமும்}

இதுவரை பார்த்தது ஒரு சாதகமான பக்கம். மறுபக்கம்,
தேர்வு இல்லாமல் நிலைமை \emph{சிக்கலாகிறது}. ஏன் என்பதைப்
பார்க்க, \cite{Hartogs1915} க்குரிய ஒரு பயனுள்ள முடிவு:

\begin{lem}[\emph{$\ZF$ இல்}]\ollabel{HartogsLemma}
எந்தக் கணம் $A$ க்கும், $\cardnless{\alpha}{A}$ ஆகும் ஒரு
வரிசையெண் $\alpha$ உள்ளது.
\end{lem}

\begin{proof}
$B \subseteq A$ மற்றும் $R \subseteq B^2$ எனில்,
\olref[ord-arithmetic][using-addition]{rankcomputation} இன்படி
$\tuple{B, R} \subseteq V_{\setrank{A}+4}$.
எனவே பிரிப்பைப் பயன்படுத்திப் பின்வருவதை எடுத்துக்கொள்க:
\[
	C = \Setabs{\tuple{B, R} \in V_{\setrank{A}+5}}{B\subseteq A
	\text{ மற்றும் $\tuple{B, R}$ ஒரு நன்கு வரிசைப்படுத்தல்}}
\]
மாற்றீட்டையும்
\olref[ordinals][ordtype]{thmOrdinalRepresentation} ஐயும்
பயன்படுத்திப் பின்வரும் கணத்தை அமைக்கவும்:
\[
	\alpha = \Setabs{\ordtype{B, R}}{\tuple{B, R} \in C}.
\]
\olref[ordinals][basic]{corordtransitiveord} இன்படி
$\alpha$ ஒரு வரிசையெண்; அது வரிசையெண்களைக் கொண்ட
மாறுநிலைக் கணம். உண்மையில் $\gamma \in \beta \in \alpha$
எனில், ஏதேனும் நன்கு வரிசைப்படுத்தலுக்கு
$\beta = \ordtype{B, R}$.
% SOURCE CORRECTION TA-STH-039: B என்பது A இன் துணைக்கணம்; R இன் துணைக்கணம் அல்ல.
இங்கு $B \subseteq A$; அப்போது ஏதேனும் $b \in B$ க்கு
\olref[ordinals][iso]{wellordinitialsegment} இன்படி
$\gamma = \ordtype{B_b, R_b}$. எனவே $\gamma \in \alpha$.

மறுதலிப்புக்கு, $f \colon \alpha \to A$ என்ற
!!a{injection} உள்ளது என்க. அப்போது,
% SOURCE CORRECTION TA-STH-040: f இன் வரையறைக்குட்பட்ட வேறுபட்ட ஆயக் குறிகளைப் பயன்படுத்துக.
\begin{align*}
	B &= \ran{f}\\
	R &= \Setabs{\tuple{f(\xi), f(\eta)} \in A \times A}{\xi \in \eta \in \alpha}.
\end{align*}
தெளிவாக $\alpha = \ordtype{B, R}$ மற்றும்
$\tuple{B, R} \in C$. எனவே $\alpha \in \alpha$;
இது முரண்.
\end{proof}

% SEGMENT OLP-0595-S02
இதிலிருந்து ஓர் ஆழமான முடிவு கிடைக்கிறது:

\begin{thm}[\emph{$\ZF$ இல்}]
பின்வரும் கூற்றுகள் சமமானவை:
\begin{enumerate}
	\item\ollabel{equivwo} நன்கு வரிசைப்படுத்தல் அடிகோளம்.
	\item\ollabel{equivcompare} எந்தக் கணங்கள் $A$, $B$
	க்கும் $\cardle{A}{B}$ அல்லது $\cardle{B}{A}$.
\end{enumerate}
\end{thm}

\begin{proof}
% SEGMENT OLP-0595-S03
\emph{\olref{equivwo} $\Rightarrow$ \olref{equivcompare}.}
$A$, $B$ ஐ நிலைப்படுத்துக. \olref{equivwo} இன்படி
$\tuple{A, R}$, $\tuple{B, S}$ என்ற நன்கு வரிசைப்படுத்தல்கள்
உள்ளன. \olref[ordinals][ordtype]{thmOrdinalRepresentation}
இன்படி $f \colon \alpha \to \tuple{A, R}$ மற்றும்
$g \colon \beta \to \tuple{B, S}$ என்ற ஒத்துருவாக்கங்கள்
உள்ளன என்க. \olref[sth][ordinals][basic]{ordinalsaresubsets}
இன்படி $\alpha \subseteq \beta$ அல்லது
$\beta \subseteq \alpha$. $\alpha \subseteq \beta$
எனில் $\comp{f^{-1}}{g} \colon A \to B$ ஓர்
!!a{injection}; எனவே $\cardle{A}{B}$. இதேபோல்
$\beta \subseteq \alpha$ எனில் $\cardle{B}{A}$.

% SEGMENT OLP-0595-S04
\emph{\olref{equivcompare} $\Rightarrow$ \olref{equivwo}.}
$A$ ஐ நிலைப்படுத்துக. \olref{HartogsLemma} இன்படி
$\cardnless{\beta}{A}$ ஆகும் ஒரு வரிசையெண் $\beta$
உள்ளது. \olref{equivcompare} இன்படி
$\cardle{A}{\beta}$. ஆகவே $f \colon A \to \beta$
என்ற ஓர் !!{injection} உள்ளது. இதன் மூலம்,
$\Setabs{\tuple{a, b} \in A \times A}{f(a) \in f(b)}$
என்ற வரிசையை வரையறுத்து $A$ இன் உறுப்புகளை நன்கு
வரிசைப்படுத்தலாம்.
\end{proof}
\noindent

% SEGMENT OLP-0595-S05
உடனடி விளைவாக, நன்கு வரிசைப்படுத்தல் தவறினால்,
சில கணங்களின் அளவுகளை உண்மையிலேயே ஒன்றோடொன்று
\emph{ஒப்பிட முடியாது}. ஆகவே அது தவறினால் முடிவிலிக்கடந்த
கண அளவெண் எண்கணிதமும் சிக்கலாகும். எடுத்துக்காட்டாக,
$A$, $B$ இரண்டும் முடிவுறாதவை எனில், $A$ மற்றும் $B$ இல்
பெரியது $M$
என்று கொண்டு வெட்டாக் கூட்டுத்தொகை மற்றும் பெருக்கற்கணம்
இரண்டுக்கும் அந்தப் பெரிய அளவே கிடைக்கும் என்ற கருத்தை
கைவிட வேண்டும்:
% SOURCE CORRECTION TA-STH-041: சமஎண்ணிக்கைச் சோதனையை உள்ளமைக்காமல் இரண்டு சரியான சமஎண்ணிக்கைச் கூற்றுகளைத் தனித்துக் கூறுக.
$\cardeq{A \disjointsum B}{M}$ மற்றும்
$\cardeq{A \times B}{M}$
(\olref[card-arithmetic][simp]{cardplustimesmax} ஐக் காண்க).
சிக்கல் எளியது: $A$, $B$ இன் அளவுகளை \emph{ஒப்பிட}
முடியாதபோது, எது பெரியது என்று கேட்பதே பொருளற்றது.

\end{document}
